\section{Contoh Terhitung}

Contoh-contoh berikut memakai angka dari bahan kuliah dan dapat diperiksa ulang
langkah demi langkah. Semua hasilnya telah dihitung ulang dengan Python.

\begin{example}
\textbf{Substitusi satu byte dengan SubBytes.} Byte $\code{23}$ mempunyai digit
atas 2 dan digit bawah 3. Nilai S-box pada baris 2 kolom 3 adalah $\code{26}$.
Dengan demikian,
\[
\text{S-box}(\code{23}) = \code{26}.
\]
Byte itu tetap berada pada posisinya di \texttt{state}; hanya nilainya yang
berubah dari $\code{23}$ menjadi $\code{26}$.
\end{example}

\begin{example}
\textbf{MixColumns satu kolom.} Ambil kolom
$[\,\code{26}\ \code{7B}\ \code{BD}\ \code{43}\,]^{T}$. Menurut
Persamaan~\eqref{eq:aes-mixcolumns},
\[
\begin{bmatrix}
\code{02} & \code{03} & \code{01} & \code{01} \\
\code{01} & \code{02} & \code{03} & \code{01} \\
\code{01} & \code{01} & \code{02} & \code{03} \\
\code{03} & \code{01} & \code{01} & \code{02}
\end{bmatrix}
\begin{bmatrix}
\code{26} \\ \code{7B} \\ \code{BD} \\ \code{43}
\end{bmatrix}
=
\begin{bmatrix}
\code{3F} \\ \code{4F} \\ \code{F9} \\ \code{2A}
\end{bmatrix}.
\]

Baris pertama dihitung sebagai berikut. Di $\text{GF}(2^8)$,
\[
\code{02} \bullet \code{26}
= x \cdot (x^5 + x^2 + x)
= x^6 + x^3 + x^2
= \code{01001100} = \code{4C}.
\]
Untuk koefisien $\code{03}$,
\[
\begin{aligned}
\code{03} \bullet \code{7B}
&= (x + 1)(x^6 + x^5 + x^4 + x^3 + x + 1) \\
&= x^7 + x^3 + x^2 + 1
= \code{10001101} = \code{8D}.
\end{aligned}
\]
Selanjutnya $\code{01} \bullet \code{BD} = \code{BD}$ dan
$\code{01} \bullet \code{43} = \code{43}$. Semua hasil antara itu di-XOR-kan:
\[
\code{4C} \oplus \code{8D} \oplus \code{BD} \oplus \code{43}
= \code{3F}.
\]
Jadi elemen pertama hasilnya adalah $\code{3F}$. Persamaan lainnya diselesaikan
dengan cara yang sama, sehingga diperoleh $\code{4F}$, $\code{F9}$, dan
$\code{2A}$.
\end{example}
